\documentclass[twocolumn]{aastex631}
\usepackage{amsmath}
\usepackage{booktabs}
\shorttitle{Optical denominator for gas-fraction versus efficiency tests}
\shortauthors{NebulaMind Research Autopilot}
\begin{document}

\title{Optical denominator for gas-fraction versus efficiency tests}
\author{NebulaMind Research Autopilot}

\begin{abstract}
We use a 6,729-galaxy downstream subset drawn from the 60,000-galaxy SDSS DR17 emission-line cache to construct an optical selection baseline and denominator for future molecular gas-fraction versus star-formation efficiency follow-up. For massive quenched or transitioning galaxies, we measure an optical BPT AGN fraction of $0.549 \pm 0.006$ and a median log H$\alpha$ luminosity proxy of 40.06 erg s$^{-1}$, which is offset by $-0.66$ dex relative to massive star-forming controls. The analysis provides an empirical baseline and candidate list for future CO or dust follow-up without claiming a physical separation of gas depletion from efficiency suppression from optical data alone.
\end{abstract}

\keywords{galaxies: evolution --- galaxies: active --- galaxies: star formation --- surveys --- methods: data analysis}
\software{Astropy, SciPy, NumPy, Matplotlib, pandas}

\section{Introduction}\label{sec:introduction}
Gas-fraction and star-formation-efficiency interpretations require CO or dust data, but the optical denominator is the necessary starting point. Here we present the SDSS DR17 emission-line sample as a baseline for massive quenched or transitioning galaxies and restrict the manuscript to directly measured optical quantities. Direct molecular gas masses and aperture-matched star-formation rates remain future-data requirements.


\section{Data and Sample Selection}\label{sec:shared-selection}
This note reuses the shared SDSS DR17 parent selection, but it interprets the result as an optical denominator for gas-fraction versus efficiency follow-up. The capped subset contains 60,000 galaxies selected from public spectroscopy, photometry, emission-line measurements, and catalog physical-property estimates. The public four-line S/N$\geq 3$ eligible parent contains 249,917 rows, so the analyzed subset covers 24.0\% of that parent. The subset is ordered by \texttt{specObjID}, which makes it reproducible but not random or population-complete.

\begin{deluxetable*}{lrrr}
\tabletypesize{\scriptsize}
\tablecaption{Shared SDSS DR17 selection cascade used before paper-specific quantities.\label{tab:selection-cascade}}
\tablehead{\colhead{Selection stage} & \colhead{Public DR17 rows} & \colhead{Cached rows} & \colhead{Retention vs. spectro-z parent (fraction)}}
\startdata
SpecObj GALAXY, 0.02<z<0.12 & 501,060 & -- & 1.000 \\
plus galSpecInfo/PhotoObj/galSpecExtra and mass/sSFR bounds & 416,554 & -- & 0.831 \\
plus galSpecLine join & 416,554 & -- & 0.831 \\
four BPT lines positive with positive errors & 373,445 & 60,000 & 0.745 \\
four BPT lines S/N$\geq 3$ & 249,917 & 60,000 & 0.499 \\
four BPT lines S/N$\geq 5$ & 176,523 & 42,446 & 0.352 \\
four BPT lines S/N$\geq 10$ & 91,768 & 22,311 & 0.183 \\
Massive quenched or transitioning subset & -- & 6,729 & -- \\
\enddata
\tablecomments{Counts are read-only public SDSS DR17 count queries plus the cached local CSV. Cached rows are shown only where the cache applies. The final row defines the specific 6,729-galaxy subset used in this optical baseline.}
\end{deluxetable*}

The paper-specific downstream selection retains 6,729 massive quenched or transitioning galaxies from the cached 60,000-galaxy parent. That branch is the denominator used for the gas-fraction versus efficiency result below, and it is a local subset rather than a separate public DR17 count query.

The four-line requirement is strongly selection dependent. In the public counts, S/N$\geq 3$ keeps 33.6\% of the $-12<\log {\rm sSFR}<-11$ parent bin but 94.9\% of the $-10<\log {\rm sSFR}<-9.5$ bin. Therefore every incidence, quenching, density, gas-denominator, or target-vector statement in this integration is conditional on the four-line emission-line selection.

Cached-versus-public marginal checks found no redshift, stellar-mass, or sSFR bin with a cached-minus-public fraction difference above 5 percentage points. The largest absolute differences were 2.03 percentage points in redshift, -1.63 percentage points in stellar mass, and -0.58 percentage points in sSFR. This is a representativeness diagnostic only; it does not make the capped cache random or complete.


\section{Measurements}\label{sec:measurements}
The row-level measurements include redshift, stellar mass, catalog specific star-formation rate, model colors, and four optical line fluxes/errors. BPT labels and all derived denominators are recomputed locally from the cached table \citep{baldwin1981,kewley2001,kauffmann2003bpt,kewley2006}. Catalog SFR/sSFR values are treated as low-redshift SDSS physical-property estimates rather than direct resolved gas or feedback measurements \citep{sdssdr17,brinchmann2004,york2000}.
Unless otherwise noted, quoted fraction uncertainties are binomial counting uncertainties from the stated sample sizes, and bracketed intervals are bootstrap confidence intervals.


\section{Optical denominator for gas-fraction versus efficiency tests}\label{sec:topic-result}
We quantify how many massive quenched or transitioning SDSS galaxies with valid emission-line measurements are available as a denominator for CO gas-fraction and depletion-time follow-up. The result is an optical baseline rather than a physical gas-depletion measurement.

The massive transition/quenched denominator contains 6,729 galaxies in the SDSS emission-line sample. Its optical BPT AGN fraction is $0.549 \pm 0.006$, and the median H$\alpha$ luminosity proxy is $\log(L_{\mathrm{H}\alpha} / \mathrm{erg~s}^{-1}) = 40.06$. The median H$\alpha$ luminosity proxy is approximately 0.66 dex lower than in massive star-forming emission-line galaxies. SDSS optical data cannot distinguish molecular-gas depletion from reduced star-formation efficiency; this paper identifies the CO follow-up denominator and optical baseline.


\begin{figure}
\centering
\includegraphics[width=\columnwidth]{../figures/fig-topic.pdf}
\caption{SDSS DR17 optical denominator/proxy diagnostic for the gas-fraction versus efficiency transition vector. The figure isolates the 6,729 massive quenched or transitioning galaxies, revealing an optical BPT AGN fraction of $0.549 \pm 0.006$.}
\label{fig:topic}
\end{figure}

\section{Interpretation and missing observables}\label{sec:missing}
This SDSS-only pilot is intentionally limited to optical quantities. A full proposal requires CO or dust-based molecular gas masses, aperture-matched SFRs, morphology, and environment labels.

Gas-fraction and depletion-time claims require CO/HI or equivalent gas masses plus aperture-matched SFRs; optical H$\alpha$ proxy values alone cannot distinguish gas depletion from low efficiency \citep{coldgass1,coldgass2,xcoldgass2017,xgass2018}.


\section{Data Availability}\label{sec:data-avail}
The SDSS DR17 spectroscopy, photometry, emission-line measurements, and catalog quantities used here are publicly available through the SDSS data releases and associated catalog papers cited below. A local subset and manifest are retained in the project repository for reproducibility.

\section{Conclusion}\label{sec:conclusion}
We have mapped the optical baseline for 6,729 massive quenched or transitioning galaxies in the SDSS emission-line sample. We find a BPT AGN fraction of $0.549 \pm 0.006$ and a median H$\alpha$ luminosity proxy of $\log(L_{\mathrm{H}\alpha} / \mathrm{erg~s}^{-1}) = 40.06$, approximately 0.66 dex lower than star-forming counterparts. While these quantities define the target selection denominator for future CO gas-fraction versus efficiency programs, direct molecular gas masses and aperture-matched star formation rates remain required to physically distinguish depletion from low efficiency.

\acknowledgments
We thank the SDSS collaboration. This manuscript uses public SDSS DR17 data only.


\begin{thebibliography}{}
\bibitem[Abdurro'uf et al.(2022)]{sdssdr17} Abdurro'uf, Accetta, K., Aerts, C., et al. 2022, ApJS, 259, 35
\bibitem[Baldwin et al.(1981)]{baldwin1981} Baldwin, J.~A., Phillips, M.~M., \& Terlevich, R. 1981, PASP, 93, 5
\bibitem[Brinchmann et al.(2004)]{brinchmann2004} Brinchmann, J., Charlot, S., White, S.~D.~M., et al. 2004, MNRAS, 351, 1151
\bibitem[Kauffmann et al.(2003a)]{kauffmann2003bpt} Kauffmann, G., Heckman, T.~M., Tremonti, C., et al. 2003a, MNRAS, 346, 1055
\bibitem[Kewley et al.(2001)]{kewley2001} Kewley, L.~J., Dopita, M.~A., Sutherland, R.~S., Heisler, C.~A., \& Trevena, J. 2001, ApJ, 556, 121
\bibitem[Kewley et al.(2006)]{kewley2006} Kewley, L.~J., Groves, B., Kauffmann, G., \& Heckman, T. 2006, MNRAS, 372, 961
\bibitem[York et al.(2000)]{york2000} York, D.~G., Adelman, J., Anderson, J.~E., Jr., et al. 2000, AJ, 120, 1579
\bibitem[Catinella et al.(2018)]{xgass2018} Catinella, B., Saintonge, A., Janowiecki, S., et al. 2018, MNRAS, 476, 875
\bibitem[Saintonge et al.(2011a)]{coldgass1} Saintonge, A., Kauffmann, G., Kramer, C., et al. 2011a, MNRAS, 415, 32
\bibitem[Saintonge et al.(2011b)]{coldgass2} Saintonge, A., Kauffmann, G., Wang, J., et al. 2011b, MNRAS, 415, 61
\bibitem[Saintonge et al.(2017)]{xcoldgass2017} Saintonge, A., Catinella, B., Tacconi, L.~J., et al. 2017, ApJS, 233, 22
\end{thebibliography}

\end{document}
